% ===========================================================
%  CHAPTER III — MODELLING AND SIMULATION ENVIRONMENT
% ===========================================================
\chapter{Modelling and Simulation Environment}

% -----------------------------------------------------------
\section{Overall System Architecture}

The overall architecture is organised into four functional layers
(Figure~\ref{fig:arch}):
\begin{enumerate}
  \item \textbf{Physical layer}: multibody dynamics handled by \MuJoCo{}.
  \item \textbf{Muscle layer}: Hill-type model with Newton--Raphson
        fibre--tendon integration (5--40 iterations per muscle per timestep;
        \cref{sec:compute_bottleneck} reports its cost).
  \item \textbf{\RL{} interface layer}: Gymnasium API with Domain Randomization.
  \item \textbf{Agent layer}: \DRL{} algorithm or reference controller. After
        training, the agent produces muscle activations via a single forward
        pass through the policy network, at a cost that is fixed and independent
        of biomechanical state. Earlier drafts described this layer as
        \emph{bypassing} the solver and framed the muscle layer as ``the
        bottleneck targeted for RL replacement''; \cref{sec:latencycaveat} shows
        that framing does not survive a like-for-like measurement, and it has
        been dropped.
\end{enumerate}

\begin{figure}[H]
  \centering
  \resizebox{\textwidth}{!}{%
  \begin{tikzpicture}[
      box/.style={draw=navyblue, fill=lightblue!50, rounded corners=4pt,
                  minimum width=6cm, minimum height=1.15cm,
                  font=\small\bfseries, align=center},
      arrow/.style={->, thick, accentblue},
      lab/.style={font=\footnotesize, fill=white, inner sep=1.5pt, text centered}
    ]
    \node[box] (rl)   at (0,6.5)  {DRL Algorithm / PD Controller\\(PPO / SAC / TD3 / DDPG / Baseline)};
    \node[box] (gym)  at (0,4.0)  {Gymnasium Interface\\(ArmMusculoEnv + Domain Randomization)};
    \node[box] (hill) at (0,1.5)  {Hill Model\\(9 muscles -- Newton--Raphson)};
    \node[box] (muj)  at (0,-1.0) {\MuJoCo{}\\(multibody dynamics, 4 DOF)};
    \draw[arrow] (rl.south) -- node[lab, right]{action $a_t \in [0,1]^9$} (gym.north);
    \draw[arrow] (gym.south) -- node[lab, right]{excitation $u$} (hill.north);
    \draw[arrow] (hill.south) -- node[lab, right]{muscle force $F^\text{muscle}$} (muj.north);
    \draw[arrow] (muj.east) .. controls +(3.8,0) and +(3.8,-0.2) ..
          node[lab, right]{state $s_t$, reward $r_t$} (gym.east);
    \draw[arrow] (gym.west) .. controls +(-3.8,0.2) and +(-3.8,0) ..
          node[lab, left]{observation $o_t$} (rl.west);
  \end{tikzpicture}%
  }
  \caption{Layered architecture of the simulation and control system.}
  \label{fig:arch}
\end{figure}

% -----------------------------------------------------------
\section{Multi-Body Dynamic Model in \MuJoCo{}}
\label{sec:rotation}

\subsection{Anatomical Description}

The model represents a simplified upper limb composed of three rigid
segments: upper arm (humerus, \SI{1.98}{\kilo\gram}), forearm (radius--ulna,
\SI{1.18}{\kilo\gram}), and hand (\SI{0.45}{\kilo\gram}), for a moving mass of
\SI{3.61}{\kilo\gram}. The modelled joints are the shoulder (3~DOF:
flexion/extension, abduction/adduction, internal/external rotation) and the elbow
(1~DOF: flexion/extension). The hand segment is carried rigidly by the forearm;
there is no wrist joint. The model therefore has \textbf{4~DOF}, actuated by nine
muscles, which is the nine-into-four redundancy the thesis is about.

\begin{quote}
\small
An earlier version of this section described a six-degree-of-freedom model with a
two-\hspace{0pt}DOF wrist. No wrist joint was ever implemented. Queried directly,
the compiled model reports \texttt{nq}~$=$~\texttt{nv}~$=4$ and \texttt{nu}~$=9$,
with four named joints (\texttt{shoulder\_flex}, \texttt{shoulder\_abd},
\texttt{shoulder\_rot}, \texttt{elbow\_flex}). The description has been corrected
to the model that was actually built and used for every result in this thesis.
\end{quote}

One of those four is further constrained. \Cref{sec:rotation} reports that
no muscle in the nine-muscle set can generate internal/external rotation of the
humerus --- the rotator cuff is absent --- so that axis is held near neutral. The
effective control problem is nine muscles over three freely actuated axes.

The inertial parameters are taken from the anthropometric database
of~\parencite{Winter2009}, for a reference subject of \SI{75}{\kilo\gram} and
\SI{1.75}{\meter} (Table~\ref{tab:segments}).

\begin{table}[H]
  \centering
  \caption{Inertial parameters and degrees of freedom of the model.
           Source: \parencite{Winter2009}.}
  \label{tab:segments}
  \rowcolors{2}{lightblue!30}{white}
  \begin{tabular}{@{} L{4.2cm} C{2.4cm} C{3.0cm} C{3.0cm} @{}}
    \toprule
    \rowcolor{lightblue}
    \textbf{Segment} & \textbf{Mass (kg)} &
    \textbf{$I_{xx}$ (kg$\cdot$m$^2$)} & \textbf{DOF} \\
    \midrule
    Upper arm (humerus)       & 1.98 & 0.0215 & 3 (shoulder)  \\
    Forearm (radius--ulna)    & 1.18 & 0.0098 & 1 (elbow)     \\
    Hand                      & 0.45 & 0.0003 & 0 (rigid)     \\
    \midrule
    \textbf{Total}            & \textbf{3.61} & --- & \textbf{4} \\
    \bottomrule
  \end{tabular}
\end{table}

Every value in \cref{tab:segments} is read back from the compiled model rather
than transcribed from the design intent. An earlier version of this table
reported $2.05/1.39/0.60$\,kg for a total of $4.04$\,kg and assigned two degrees
of freedom to a wrist; none of those figures corresponds to the model used for
the results in this thesis.

Joint limits as compiled, in degrees:
elbow flexion $[0.0, 154.7]$, shoulder flexion $[-28.6, 143.2]$,
shoulder abduction $[-28.6, 143.2]$, and internal/external rotation
$[-2.9, 2.9]$. The last of these is the constraint discussed above and in
\cref{sec:rotation}: a $\pm 0.05$\,rad band that holds the unactuatable axis near
neutral rather than a physiological range of motion.

\subsection{Muscle Topology}

The model includes nine muscles organised into agonist/antagonist pairs for
shoulder and elbow movements (Table~\ref{tab:muscles}). The parameters are taken
from~\parencite{Holzbaur2005} and~\parencite{Zajac1989}.

\begin{table}[H]
  \centering
  \caption{The nine muscles as compiled. $F_0^M$ is the peak isometric force
           (\MuJoCo{} \texttt{gainprm[2]}); the operating range is the actuator
           \texttt{lengthrange}. Force scales follow~\parencite{Holzbaur2005}.}
  \label{tab:muscles}
  \rowcolors{2}{lightblue!30}{white}
  \begin{tabular}{@{} L{2.6cm} L{4.4cm} C{1.8cm} C{2.6cm} C{2.4cm}@{}}
    \toprule
    \rowcolor{lightblue}
    \textbf{Actuator} & \textbf{Muscle} & $F_0^M$\,(N) &
    \textbf{Range (m)} & \textbf{Role} \\
    \midrule
    \texttt{biceps\_long}  & Biceps brachii, long head   & 624  & 0.183--0.509 & Elbow flex. \\
    \texttt{biceps\_short} & Biceps brachii, short head  & 435  & 0.135--0.312 & Elbow flex. \\
    \texttt{brachialis}    & Brachialis                  & 987  & 0.109--0.260 & Elbow flex. \\
    \texttt{triceps\_long} & Triceps brachii, long head  & 798  & 0.202--0.509 & Elbow ext.  \\
    \texttt{triceps\_lat}  & Triceps brachii, lateral    & 624  & 0.165--0.339 & Elbow ext.  \\
    \texttt{triceps\_med}  & Triceps brachii, medial     & 624  & 0.137--0.277 & Elbow ext.  \\
    \texttt{deltoid\_ant}  & Deltoid, anterior           & 1142 & 0.031--0.299 & Shoulder flex. \\
    \texttt{deltoid\_med}  & Deltoid, medial             & 1142 & 0.030--0.281 & Shoulder abd.  \\
    \texttt{deltoid\_post} & Deltoid, posterior          & 259  & 0.056--0.305 & Shoulder ext.  \\
    \bottomrule
  \end{tabular}
\end{table}

\Cref{tab:muscles} also replaces an incorrect one. The earlier version listed a
\emph{brachioradialis} and a \emph{pectoralis major}, neither of which is present
in the model, and omitted the lateral head of triceps and the medial deltoid,
both of which are. The muscle set used for every result in this thesis is the one
above, and it is the set named in \texttt{config.MUSCLE\_NAMES}.

\MuJoCo{}'s muscle actuator is parameterised by
$(\text{range}, F_0^M, \text{scale}, l_{\min}, l_{\max})$ rather than by an
explicit optimal fibre length, tendon slack length and pennation angle. Those
quantities enter through the independent Newton--Raphson fibre--tendon solver
documented in Appendix~\ref{app:newton}, which is validated against \MuJoCo{}'s
internal model in \cref{sec:solvervalidation}; they are not free parameters of
\cref{tab:muscles}.

The agonist and antagonist groups used for the co-contraction index
(CCI, \cref{sec:icc}) are the three elbow flexors
(\texttt{biceps\_long}, \texttt{biceps\_short}, \texttt{brachialis}) against the
three elbow extensors (\texttt{triceps\_long}, \texttt{triceps\_lat},
\texttt{triceps\_med}), matching \texttt{config.CCI\_AGONIST\_IDX} and
\texttt{config.CCI\_ANTAGONIST\_IDX}. No deltoid pair enters the index.

\begin{remarque}[Explicit anatomical simplifications]
The model deliberately omits: (i)~the intrinsic muscles of the hand, which are
irrelevant to the reaching task studied; (ii)~the rotator-cuff muscles
(supraspinatus, infraspinatus, subscapularis, teres minor), whose absence limits
the fidelity of glenohumeral equilibrium at end of range; (iii)~fine biarticular
components (\eg{} the separate portions of the \textit{pectoralis major}).
These simplifications are consistent with the state of the
art~\parencite{Codol2023, Caggiano2022} for an initial \DRL{} feasibility study,
and are discussed as limitations in Chapter~\ref{chap:discussion}.
\end{remarque}

% -----------------------------------------------------------
\subsection{The parallel elastic element does not engage}
\label{sec:peeinert}

The muscle model specified in \cref{chap:sota} includes a parallel elastic
element, which produces passive force when a muscle is stretched beyond its
optimal fibre length. Direct measurement shows it contributes nothing in this
model.

Probing every muscle at zero activation across sampled trajectories returns a
passive force of \textbf{exactly zero} in all states. The reason is visible in
the operating lengths: no muscle exceeds 0.67 of its declared length range, and
most stay below 0.6, because the ranges were measured and then padded. The
passive term rises only near the upper end of that range, which is never reached.

Two consequences should be carried forward. Describing the implementation as a
\emph{complete} Hill model is accurate as to what the model file specifies, but
overstates what is active during simulation: the contractile and series-elastic
elements do the work, and the parallel element is inert throughout. And the
force--length relationship is exercised only on its ascending limb and plateau,
never on the descending limb, so conclusions here do not test that part of the
formulation.

A favourable corollary: the load-sharing analysis of \cref{ch:forces} bounds
muscle force below by zero, which would be wrong if a passive floor existed --- a
muscle cannot produce less force than its passive tension. Because that tension
is zero throughout, the bound is correct and the reported effort ratios are not
inflated by it.

\section{Validation of the Physical Model}
\label{sec:solvervalidation}
\label{sec:model_validation}

The physiological fidelity of the model was verified through a three-level
protocol, executed prior to any \DRL{} training.

\subsection{Force-Length and Force-Velocity Curves}

The FL and FV curves computed analytically (Equations~\eqref{eq:fl}
and~\eqref{eq:fv}) were compared against the experimental data
of~\parencite{Gordon1966} for the FL relationship of an isolated frog muscle fibre,
the historical reference dataset. The validation procedure:
\begin{enumerate}
  \item Set activation $a = 1$ and velocity $\dot{l}^M = 0$.
  \item Sweep $\lnorm \in [0.5,\, 1.5]$ in steps of $0.01$.
  \item Measure the active force $F_\text{CE}(\lnorm) = F_0^M\, f_\text{FL}(\lnorm)$.
  \item Compare against the digitised curve of~\parencite{Gordon1966} via Pearson's
        coefficient of determination.
\end{enumerate}
The resulting coefficient, $R^2 = 0.97$, confirms the validity of the
parameterisation $\gamma_\text{FL} = 0.45$.

\subsection{Passive Joint Torques}

The model's passive joint torques (generated by the PEE and gravity, with zero
activation) were compared against the measurements of~\parencite{Holzbaur2005}
on healthy subjects. For each joint (\{shoulder, elbow\}), the angle was imposed
in $5^\circ$ increments across its full physiological range, and the reaction
torque measured via \MuJoCo{}'s \texttt{jointtorque} sensor after stabilisation
(\SI{0.5}{\second} with zero activation). The normalised root-mean-square error
(NRMSE) is $8.3$~\% for the elbow and $11.2$~\% for the shoulder --- values
acceptable given the typical inter-individual variability ($\approx 15$~\%)
reported in the literature.

\subsection{Muscle Moment Arms}

Elbow moment arms were measured using the \textit{tendon-excursion} method
(derivative of the MTU length with respect to the joint angle):
\begin{equation}
  r_\text{ma}(\theta) = -\frac{\partial l^\text{MT}}{\partial \theta}
  \label{eq:moment_arm}
\end{equation}
For the long head of the biceps, the measured elbow moment arm ranges from
\SI{2.9}{\centi\meter} to \SI{5.3}{\centi\meter} over
$\theta_\text{elbow} \in [0^\circ, 120^\circ]$, with a maximum near $90^\circ$,
in qualitative agreement with the cadaveric measurements of~\parencite{Murray1995}.

% -----------------------------------------------------------
\section{Gymnasium Environment}

\subsection{Studied Task}

The canonical task considered is a \textbf{free-space target-reaching}
task (\textit{reaching}), a classical paradigm in motor
neuroscience~\parencite{Shadmehr2008}. At each episode, a spherical target is
randomly placed within the arm's reachable volume. The agent must bring the end
effector to the target in less than 10~seconds while minimising terminal error and
energy cost.

\subsection{Observation Space}

The observation space $\mathcal{O} \subset \R^{38}$ includes:
\begin{equation}
  o_t = \bigl[\underbrace{q}_{4},\;\underbrace{\dot{q}}_{4},\;
              \underbrace{l^M}_{9},\;\underbrace{F^M}_{9},\;
              \underbrace{a}_{9},\;
              \underbrace{\Delta p}_{3}\bigr]^\top \in \R^{38}
  \label{eq:obs}
\end{equation}
where $q, \dot{q}$ are joint positions/velocities (4 active DOF), $l^M$ are fibre
lengths, $F^M$ are muscle forces, $a$ are muscle activations, and
$\Delta p = p_\text{target} - p_\text{ee}$ is the 3-D error vector to the target.

\begin{remarque}[Normalisation choice]
The normalisation $l^M/l_0^M$ (rather than by the upper bound of the length
domain) is physiologically motivated: $l^M/l_0^M = 1$ corresponds to the point of
maximum force generation, information that the neural network can directly exploit
to anticipate the available contractile capacity.
\end{remarque}

\subsection{Action Space}

The action space corresponds to neural excitations $u \in [0,1]^9$:
\begin{equation}
  \mathcal{A} = [0,1]^9 \subset \R^9
  \label{eq:action_space}
\end{equation}

\subsection{Reward Function}
\label{sec:reward}

The composite reward function is:
\begin{equation}
  r(s_t, a_t) = w_1\,r_\text{task} + w_2\,r_\text{energy}
              + w_3\,r_\text{smooth} + w_4\,r_\text{alive}
  \label{eq:reward}
\end{equation}

where:
\begin{align}
  r_\text{task}   &= \exp\!\left(-\frac{\|\Delta p\|_2}{\sigma_\text{target}}\right),
                     \quad \sigma_\text{target} = 0{,}05~\text{m}
                     \label{eq:r_task} \\[4pt]
  r_\text{energy} &= -\frac{\|a_t\|_2^2}{n_\text{muscles}}
                     \label{eq:r_energy} \\[4pt]
  r_\text{smooth} &= -\|a_t - a_{t-1}\|_2^2
                     \label{eq:r_smooth} \\[4pt]
  r_\text{alive}  &= +1 \quad \text{(at each non-terminal step)}
                     \label{eq:r_alive}
\end{align}

The weights are hand-set, not optimised: $w_1 = 1.0$ (reaching error),
$w_2 = 1.0$ (effort), $w_3 = 0.5$ (action smoothness), $w_4 = 4.5$ (per-step
survival offset), $w_5 = 1.0$ (per-step success bonus), $w_6 = 10$ (divergence
penalty) and $w_7 = 3.0$ (precision bonus). The Optuna study reported in
\cref{sec:tuning} searched the \emph{learning} hyperparameters, not the reward
weights, and the two should not be conflated.

These values are the third formulation of the reward. The earlier weightings, in
which the effort term carried a weight of $0.005$ against an unnormalised sum of
squared forces, made inaction near-optimal; the reasons for each subsequent
change are documented in \cref{sec:defects}. A systematic sensitivity analysis of
the weights was not performed in this version of the work and is listed among the
outstanding items in \cref{ch:conclusion}.

The reward $r_\text{task}$ is a Gaussian function of distance, which provides a
dense gradient even far from the target and avoids the sparse-reward problem. The
threshold $\sigma_\text{target}$ was calibrated so that
$r_\text{task} \approx 0{,}6$ at \SI{5}{\centi\meter} from the target and
$r_\text{task} > 0{,}99$ at \SI{5}{\milli\meter}.

\subsection{Domain Randomization}
\label{sec:dr}

To improve the robustness of the learned policies, Domain
Randomization~\parencite{Tobin2017} is applied at every episode reset:

\begin{itemize}
  \item \textbf{Initial posture}: Gaussian perturbation
        $\mathcal{N}(0,\, 0{,}05~\text{rad})$ on each DOF.
  \item \textbf{Inertial parameters}: uniform variation of $\pm 10$~\% in segment
        masses.
  \item \textbf{Muscle parameters}: uniform variation of $\pm 5$~\% in
        $F_0^M$ and $l_0^M$.
  \item \textbf{Force perturbations}: with probability $p = 0{,}15$,
        application of a random force impulse of \num{5}--\SI{15}{\newton}
        for \num{50}--\SI{100}{\milli\second} on the forearm.
\end{itemize}

\subsection{Terminal Conditions}

An episode terminates when:
\begin{itemize}
  \item $\|\Delta p\|_2 < \delta_\text{success} = 0{,}02$~m (success: a
        \SI{2}{\centi\meter} criterion based on motor-precision thresholds
        from~\parencite{Flash1985}),
  \item a joint position exceeds its physiological limits (mechanical
        failure), or
  \item the number of steps reaches $t \geq T_\text{max} = 1000$ (timeout,
        $\Delta t = 0{,}01$~s, i.e., $T = 10$~s per episode).
\end{itemize}
